\begin{problem}[31.2]
  Find all prime ideals and all maximal ideals of $\Z_{12}$.
\end{problem}

\begin{solution}
  The prime ideals of $\Z_{12}$ are given by $p_1^{k_1} \cdots p_n^{k_n}$, where each $p_i$ is a prime factor of 12. Thus, for 12, we have $12 = 2^2 \cdot 3$. Thus, the prime ideals for $\Z_{12}$ are given by $\braket{2} \aand \braket{3}$. Since we're working with $\Z_{12}$, the prime ideals are also maximal ideals. Thus, the maximal ideals of $\Z_{12}$ are also given by $\braket{2}$ and $\braket{3}$.
\end{solution}

\begin{problem}[31.4]
  Find all prime ideals and all maximal ideals of $\Z_2 \times \Z_4$.
\end{problem}

\begin{solution}
  The prime ideals of $\Z_2 \times \Z_4$ are given by $P \times \Z_4$ or $\Z_2 \times Q$, where $P$ is a prime ideal of $\Z_2$ and $Q$ is a prime ideal of $\Z_4$. The only prime ideal of $\Z_2$ is $\braket{0}$, and the prime ideals of $\Z_4$ is $\braket{2}$. Thus, the prime ideals of $\Z_2 \times \Z_4$ are given by $\braket{0} \times \Z_4$  and $\Z_2 \times \braket{2}$. The maximal ideals of $\Z_2 \times \Z_4$ are given by $M \times \Z_4$ or $\Z_2 \times N$, where $M$ is a maximal ideal of $\Z_2$ and $N$ is a maximal ideal of $\Z_4$. The only maximal ideal of $\Z_2$ is $\braket{0}$, and the maximal ideals of $\Z_4$ is $\braket{2}$. Thus, the maximal ideals of $\Z_2 \times \Z_4$ are given by $\braket{0} \times \Z_4$  and $\Z_2 \times \braket{2}$.
\end{solution}

\begin{problem}[31.6]
  Find all $c \in \Z_3$ such that $\Z_3[x]/\braket{x^3 + x^2 + c}$ is a field.
\end{problem}

\begin{solution}
  In order for $\Z_3[x]/\braket{f_c(x)}$ to be a field, where $f_c(x) = x^3 + x^2 + c$, we must have $f_c(x)$ be an irreducible polynomial. We check each of the values $c = 0, 1, 2$. For $c = 0$, we have $f_0(x) = x^3 + x^2$, which is reducible, since it has a root at $x = 0$. For $c = 1$, we have $f_0(x) = x^3 + x^2 + 1$, which is reducible, since it it has a root at $x = 1$. For $c = 2$, we have $f_2(x) = x^3 + x^2 + 2$, which is irreducible, since it doesn't have any roots. Therefore, for $\Z_3[x]/\braket{x^3 + x^2 + c}$ to be a field, we must have $c = 2$.
\end{solution}

\begin{problem}[31.8]
  Find all $c \in \Z_5$ such that $\Z_5[x]/\braket{x^2 + x + c}$ is a field.
\end{problem}

\begin{solution}
  Again, just like problem 31.6, we check for what values of $c \in \Z_5$ give us $f_c(x) = x^2 + x + c$ an irreducible polynomial. Those values of $c$ are 1 and 2, since $f_1(x)$ and $f_2(x)$ don't have any roots.
\end{solution}

\begin{problem}[31.16]
  Find a prime ideal of $\Z \times \Z$ that is not maximal.
\end{problem}

\begin{solution}
  We construct the ideal $I = p\Z \times \Z$, where $p = 2$ is a prime number. We now check that it is an ideal of $\Z \times \Z$.

  Clearly, $I$ is nonempty since $(2,0) \in I$. Let $(a_1,b_1), (a_2,b_2) \in I$. Then $a_1, a_2 \in 2\Z$ and $b_1, b_2 \in \Z$. Hence
  \[%
    (a_1,b_1) - (a_2,b_2) = (a_1 - a_2,\; b_1 - b_2)
  ,\]%
  and since $2\Z$ and $\Z$ are ideals of $\Z$, we have $a_1 - a_2 \in 2\Z$ and $b_1 - b_2 \in \Z$. Thus $I$ is closed under subtraction.

  Now let $(r,s) \in \Z \times \Z$ and $(a,b) \in I$. Then $a \in 2\Z$, so $ra \in 2\Z$, and clearly $sb \in \Z$. Hence
  \[%
    (r,s)(a,b) = (ra,\; sb) \in I
  .\]%
  Thus $I$ is an ideal of $\Z \times \Z$.

  Next, we show that $I$ is prime. Consider the quotient ring
  \[%
    (\Z \times \Z)/I
  .\]%
  Define a map
  \[%
    \phi : \Z \times \Z \to \Z/2\Z, \quad \phi(a,b) = \bar{a}
  .\]%
  This is a surjective ring homomorphism, and
  \[%
    \ker \phi = 2\Z \times \Z = I
  .\]%
  By the First Isomorphism Theorem,
  \[%
    (\Z \times \Z)/I \cong \Z/2\Z
  ,\]%
  which is a field, hence an integral domain. Therefore $I$ is a prime ideal.

  Finally, $I$ is not maximal. Indeed, consider the ideal
  \[%
    J = 2\Z \times 2\Z
  .\]%
  Then
  \[%
    I = 2\Z \times \Z \subsetneq 2\Z \times 2\Z \subsetneq \Z \times \Z
  ,\]%
  so there exists a proper ideal strictly between $I$ and the whole ring. Hence $I$ is not maximal.

  Therefore $I = 2\Z \times \Z$ is a prime ideal of $\Z \times \Z$ that is not maximal.
\end{solution}

\begin{problem}[31.18]
  Is $\Q[x]/\braket{x^2 - 5x + 6}$ is a field? Why?
\end{problem}

\begin{solution}
  The quotient $\Q[x]/\braket{x^2 - 5x + 6}$ isn't a field since we can factor the polynomial into $(x - 2)(x - 3)$. Since the polynomial is reducible, the ideal $\braket{x^2 - 5x + 6}$ isn't maximal, meaning the quotient isn't a field.
\end{solution}

\begin{problem}[31.24]
  Give an example of an ideal in $\Q[x, y]$ that is not a principal ideal. Conclude that if $R$ is an integral domain with the property that every ideal in $R$ is principal, it does not follow that every ideal in $R[x]$ is a principal ideal.
\end{problem}

\begin{solution}
  Consider the ideal
  \[%
    I = \braket{x, y} \subset \Q[x,y].
  \]%
  We claim that $I$ is not principal. We prove this via contradiction. Assume that $I = \braket{f(x,y)}$ for some $f \in \Q[x,y]$. Since $x, y \in I$, we must have
  \[%
    f \mid x \aand f \mid y
  \]%
  Thus $f$ is a common divisor of $x$ and $y$. However, $x$ and $y$ are irreducible and not associates in the UFD $\Q[x,y]$, so their only common divisors are units. Hence $f$ must be a unit, which implies
  \[%
    I = \Q[x,y]
  ,\]%
  a contradiction since $1 \notin \braket{x,y}$. Therefore $\braket{x,y}$ is not a principal ideal.

  Now let $R$ be an integral domain such that every ideal in $R$ is principal. Taking $R = \Q$, which is a field (hence a PID), we have
  \[%
    R[x][y] = \Q[x,y]
  .\]%
  Since $\Q[x,y]$ contains the non-principal ideal $\braket{x,y}$, it follows that even if every ideal in $R$ is principal, it does not follow that every ideal in $R[x]$ is principal.
\end{solution}

\begin{problem}[31.26]
  Let $R$ be a finite commutative ring with unity. Show that every prime ideal in $R$ is a maximal ideal.
\end{problem}

\begin{solution}
  Let $P$ be a prime ideal in $R$. Then the quotient ring $R/P$ is an integral domain. Since $R$ is finite and $P$ is an ideal, the quotient $R/P$ is also finite. Thus $R/P$ is a finite integral domain. We claim that every finite integral domain is a field. Let $\bar{a} \in R/P$ with $\bar{a} \neq 0$. Consider the map
  \[%
    \phi : R/P \to R/P, \quad \phi(\bar{x}) = \bar{a}\bar{x}
  .\]%
  If $\phi(\bar{x}) = \phi(\bar{y})$, then
  \[%
    \bar{a}\bar{x} = \bar{a}\bar{y} \implies \bar{a}(\bar{x} - \bar{y}) = 0
  .\]%
  Since $R/P$ is an integral domain and $\bar{a} \neq 0$, it follows that $\bar{x} = \bar{y}$. Thus $\phi$ is injective. Because $R/P$ is finite, $\phi$ is also surjective. Hence there exists $\bar{x} \in R/P$ such that
  \[%
    \bar{a}\bar{x} = 1
  \]%
  Therefore every nonzero element of $R/P$ is invertible, so $R/P$ is a field.

  Since $R/P$ is a field, $P$ is a maximal ideal. Thus every prime ideal in $R$ is maximal.
\end{solution}

\begin{problem}[31.32]
  Prove that if $F$ is a field, every proper nontrivial prime ideal of $F[x]$ is maximal.
\end{problem}

\begin{solution}
  Let $P$ be a proper nontrivial prime ideal of $F[x]$, where $F$ is a field. Since $F[x]$ is a principal ideal domain, there exists a nonzero polynomial $f(x) \in F[x]$ such that
  \[%
    P = \braket{f(x)}
  .\]%
  Because $P$ is prime, $\braket{f(x)}$ is a prime ideal. In a PID, a principal ideal is prime if and only if its generator is irreducible. Thus $f(x)$ is irreducible in $F[x]$. Consider the quotient ring
  \[%
    F[x]/\braket{f(x)}
  .\]%
  Since $f(x)$ is irreducible over the field $F$, the ideal $\braket{f(x)}$ is maximal, and hence the quotient is a field.

  Therefore $P$ is a maximal ideal. Thus every proper nontrivial prime ideal of $F[x]$ is maximal.
\end{solution}

\begin{problem}[31.36]
  If $A$ and $B$ are ideals of a ring $R$, the \textit{sum} $A + B$ of $A$ and $B$ is defined by
  \[%
    A + B = \{a + b~\mid~a \in A, b \in B\}%
  .\]%
  \begin{enumerate}
    \item Show that $A + B$ is an ideal.
    \item Show that $A \subseteq A + B$ and $B \subseteq A + B$.
  \end{enumerate}
\end{problem}

\begin{solution}[(i)]
  First, $0 = 0 + 0 \in A + B$ since $0 \in A$ and $0 \in B$. Let $x, y \in A + B$. Then there exist $a_1, a_2 \in A$ and $b_1, b_2 \in B$ such that
  \[%
    x = a_1 + b_1, \quad y = a_2 + b_2
  .\]%
  Then
  \[%
    x - y = (a_1 - a_2) + (b_1 - b_2)
  .\]%
  Since $A$ and $B$ are ideals, $a_1 - a_2 \in A$ and $b_1 - b_2 \in B$. Hence $x - y \in A + B$. Thus $A + B$ is closed under subtraction.

  Now let $r \in R$ and $x = a + b \in A + B$. Then
  \[%
    rx = ra + rb
  .\]%
  Since $A$ and $B$ are ideals, $ra \in A$ and $rb \in B$, so $rx \in A + B$. Similarly, $xr \in A + B$ if $R$ is not assumed commutative.

  Therefore $A + B$ is an ideal of $R$.
\end{solution}

\begin{solution}[(ii)]
  Since for any $a \in A$ we have $a = a + 0$ with $0 \in B$, it follows that $a \in A + B$. Thus $A \subseteq A + B$.

  Similarly, for any $b \in B$ we have $b = 0 + b$ with $0 \in A$, so $b \in A + B$. Hence $B \subseteq A + B$.
\end{solution}
